% Options for packages loaded elsewhere
\PassOptionsToPackage{unicode}{hyperref}
\PassOptionsToPackage{hyphens}{url}
%
\documentclass[
]{article}
\usepackage{amsmath,amssymb}
\usepackage{lmodern}
\usepackage{iftex}
\ifPDFTeX
  \usepackage[T1]{fontenc}
  \usepackage[utf8]{inputenc}
  \usepackage{textcomp} % provide euro and other symbols
\else % if luatex or xetex
  \usepackage{unicode-math}
  \defaultfontfeatures{Scale=MatchLowercase}
  \defaultfontfeatures[\rmfamily]{Ligatures=TeX,Scale=1}
\fi
% Use upquote if available, for straight quotes in verbatim environments
\IfFileExists{upquote.sty}{\usepackage{upquote}}{}
\IfFileExists{microtype.sty}{% use microtype if available
  \usepackage[]{microtype}
  \UseMicrotypeSet[protrusion]{basicmath} % disable protrusion for tt fonts
}{}
\makeatletter
\@ifundefined{KOMAClassName}{% if non-KOMA class
  \IfFileExists{parskip.sty}{%
    \usepackage{parskip}
  }{% else
    \setlength{\parindent}{0pt}
    \setlength{\parskip}{6pt plus 2pt minus 1pt}}
}{% if KOMA class
  \KOMAoptions{parskip=half}}
\makeatother
\usepackage{xcolor}
\setlength{\emergencystretch}{3em} % prevent overfull lines
\providecommand{\tightlist}{%
  \setlength{\itemsep}{0pt}\setlength{\parskip}{0pt}}
\setcounter{secnumdepth}{-\maxdimen} % remove section numbering
\ifLuaTeX
  \usepackage{selnolig}  % disable illegal ligatures
\fi
\IfFileExists{bookmark.sty}{\usepackage{bookmark}}{\usepackage{hyperref}}
\IfFileExists{xurl.sty}{\usepackage{xurl}}{} % add URL line breaks if available
\urlstyle{same} % disable monospaced font for URLs
\hypersetup{
  hidelinks,
  pdfcreator={LaTeX via pandoc}}

\author{Dietmar Gerald Schrausser}
\date{09.05.2026}


\begin{document}

\hypertarget{foliumblat-iiiiv}{%
\section{Foliu{[}m{]}/Blat IIIIv}\label{foliumblat-iiiiv}}

A \emph{contemporary} English translation is presented from comparing
the existing Latin, Early New High German (ENHG) and English
(translation) texts for better comprehensibility. Especially since the
ENHG text is difficult to understand compared to modern High German, as
a \emph{profound} knowledge of German (as a mother tongue) as well as
already aknowledge of various \emph{dialects} is required.

\hypertarget{on-the-work-of-the-fifth-day}{%
\subsection{On the Work of the Fifth
Day}\label{on-the-work-of-the-fifth-day}}

\begin{quote}
\textbf{Q}Uinto die dixit deus: producant aque reptile anime viuentis.
{[}et{]} volatile super terra{[}m{]} sub firmamento celi.
Cerauitq{[}ue{]} de{[}us{]} cete grandia: {[}et{]} omnem animam viuentem
atq{[}ue{]} motabilem quas produxera{[}n{]}t aque in species suas.
{[}et{]} om{[}n{]}e volatile sec{[}un{]}d{[}u{]}m genus suum.
Uide{[}n{]}s quod esset bonum benedixit eis dicens. Crescite {[}et{]}
multiplicamini {[}et{]} replete aquas maris: auesq{[}ue{]}
multiplicentur super terram.
\end{quote}

So on that day God adorned the air and the water with creatures, giving
winged ones to the air and swimming ones to the water, which are called
\emph{reptiles}\footnote{`reptilia', only in the Latin version.}, which
scamper about with a certain impetuosity. Indeed, one finds large whales
and aquatic creatures in the sea, of monstrous form, larger than land
animals due to their abundance of moisture\footnote{`habundantia
  humoris' and `uberflüssigkeit irer feuchtigkeit'.}.\footnote{from the
  `Historia Naturalis' (Plinius Secundus,
  \href{https://gallica.bnf.fr/ark:/12148/btv1b550140045/}{1250}, Book
  9, Chapter 2).} It is said that whatever arises in any part of nature
also exists in the sea.\footnote{from the `Historia Naturalis' (Plinius
  Secundus,
  \href{https://gallica.bnf.fr/ark:/12148/btv1b550140045/}{1250}, Book
  9, Chapter 1).}

Now to the following things, which are evident with regard to the origin
of animals: \emph{After} (i) plants come those (ii) beings that feel and
move; although even to plants the Pythagoreans ascribe a dull\footnote{`stupidum
  sensum' and `unbrüfende empfindlichkeit'.} sense. Such creatures,
which clearly exhibit movement and sensation, are divided by
\textbf{Moses} and in the \textbf{\emph{Timaeus}} into flying animals,
aquatic animals, and terrestrial animals.\footnote{Drawing a parallel
  between the biblical creation account (Moses) and Plato's cosmology in
  Timaeus (\href{https://doi.org/10.3931/e-rara-24349}{1588}) was a
  common theme in Renaissance Neoplatonism or Scholasticism.}

Now

let us turn to \emph{Moses}, who, after speaking about heavenly things,
mentioned the earthly animals in proper order, which live either in the
water, on the earth, or in the air.\footnote{Likely from the `Hexameron'
  (Ambrosius,
  \href{https://mdz-nbn-resolving.de/details:bsb00046506}{12th c.}).}

\begin{quote}
Si tamen inhabitare aerem volucres dici possunt.\footnote{`If one can
  even say that birds inhabit the air,' a passage from the `Heptaplus'
  (Mirandola \& Mirandola,
  \href{https://doi.org/10.3931/e-rara-61217}{1601}, Buch I, Kapitel 5),
  it questions the nature of birds and their relationship to the
  \emph{air} as a habitat in comparison to the \emph{waters} of the
  biblical creation story.}
\end{quote}

\textbf{\emph{Let us leave this discussion here}}. In what way do the
bodies of animals\footnote{`corpora animalium' and `lieb der thier'.}
arise from the elements, what are the basic reasons\footnote{`seminarie
  rationes' and `besamunge{[}n{]}'.} that God has implanted in the
nature of things, is the life of irrational animals also brought forth
from the womb of matter, or does all life rather originate from a divine
principle, as \textbf{Plotinus} most emphatically asserts? A statement
which the prophet seems to support in the following passage, after he
said:\footnote{From Ficino's commentary on Plotinus' `Enneads'
  (\href{https://doi.org/10.3931/e-rara-120889}{1580}).}

\begin{quote}
producant aque reptile anime viuentis.\footnote{'May the waters bring
  forth living, crawling creatures.}
\end{quote}

Later the following was added:

\begin{quote}
creauit deus omne anima{[}m{]} viuente{[}m{]}.\footnote{`God created
  every living soul'.}
\end{quote}

It should be noted that the \emph{water} brought these things forth at
God's (i) command, and God himself then created them. Where God's (ii)
action is mentioned, it then says

\begin{quote}
creauit deus \emph{anima{[}m{]}} viuente{[}m{]}, \footnote{`God created
  a living soul'.}
\end{quote}

however, where only \emph{water} is mentioned, instead of

\begin{quote}
\emph{a{[}n{]}imam}\footnote{`one soul'.}
\end{quote}

of

\begin{quote}
reptile \emph{a{[}n{]}i{[}m{]}e}\footnote{`Crawling creatures'. Both,
  \emph{anima} \emph{and} \emph{animus} mean \emph{soul}, but describe
  different aspects of existence: \textbf{\emph{Anima}}, the life
  principle, \textbf{\emph{Animus}}, the rational mind, c.f. Jung's
  `inner self' (s. Jung,
  \href{https://books.google.com/books?id=cX3VEAAAQBAJ}{2023}).}
\end{quote}

that is, a \emph{vehicle} for the soul provided by the water.\footnote{This
  text is a passage from `Genesis ad litteram' (Augustinus,
  \href{https://mdz-nbn-resolving.de/details:bsb00046071}{1961}, Book 3,
  Chapter 13). It deals with the linguistic and theological nuances of
  the creation of aquatic life in Genesis 1:20-21 (c.f. Jiménez de
  Cisneros, \href{https://doi.org/10.3931/e-rara-46695}{1517}).}

As Moses mentions \emph{three} animal species on Earth in the following
day of creation, the most numerous and largest animals are in the sea of
\hspace{0pt}\hspace{0pt}the Indians, including \emph{large
fish}\footnote{`balene'.} the size of four \emph{Iugera}\footnote{1
  Iugerum = 0.25 Hectares, more likely a school of fish.}, and
\emph{whales}\footnote{`bellue' and `wu{[}n{]}der thier'.} can also be
seen in the sea where the sun takes its course\footnote{The tropics.}.
Here, they are carried up from the depths of the sea by enormous waves
until they become visible to humans. Then whirlwinds rage, the rain
comes, followed by storms; tumbling down from the mountain ridges, they
churn the seas from the bottom and tear marine animals from the depths
with the waves. Among the birds, the largest, almost the largest of all
animals, are the African or Ethiopian ostriches\footnote{`generis
  strucio cameli'.}, which tower over a rider on a horse and surpass
them in speed.\footnote{from the `Historia Naturalis' (Plinius Secundus,
  \href{https://gallica.bnf.fr/ark:/12148/btv1b550140045/}{1250}, Book
  9, 10).} Many more wonderful things about the nature of birds and fish
are experienced daily in many different places\footnote{`Multo
  mirabilius de naturis auium {[}et{]} piscium ratio experiendi quotidie
  in varijs locis datur', Roger Bacon
  (\href{https://doi.org/10.3931/e-rara-16491}{1733}, Pars VI: Scientia
  Experimentalis).}.

\hypertarget{references}{%
\subsection{References}\label{references}}

Ambrosius, A. (12th c.) \emph{Ambrosii Hexaemeron Libri VI}. München:
Bayerische Staatsbibliothek.
\url{https://mdz-nbn-resolving.de/details:bsb00046506}

Augustinus, A. (1961). \emph{Über Den Wortlaut Der Genesis DE GENESI AD
LITTERAM LIBRI DUODECIM: Der Große Genesiskommentar in Zwölf Büchern Zum
Erstenmal in Dt. Sprache von Carl Johann Perl}. Edited by Perl, C. J.
Paderborn: Schöningh.
\url{https://mdz-nbn-resolving.de/details:bsb00046071}

Bacon, R. (1733). \emph{Fratris Rogeri Bacon \ldots{} opus maius
\ldots{} ex MS. codice Dubliniensi}. Edited by Jebb, S. Cum aliis
quibusdam collato, nunc primum edidit. Londini: typis Gulielmi Bowyer.
\url{https://doi.org/10.3931/e-rara-16491}

Jiménez de Cisneros, F. (1517). \emph{Biblia Polyglotta Complutensis}.
Complutum: Arnaldo Guillén de Brocar.
\url{https://doi.org/10.3931/e-rara-46695}

Jung, C. G. (2023). \emph{Collected Works of C. G. Jung: The First
Complete English Edition of the Works of C. G. Jung}. United Kingdom:
Routledge. \url{https://books.google.com/books?id=cX3VEAAAQBAJ}

Mirandola, G., \& Mirandola, G. F. (1601). \emph{Ioannis Pici,
Mirandulae \ldots{} opera quae extant omnia \ldots{}} Basileae: per
Sebastianum Henricpetri. \url{https://doi.org/10.3931/e-rara-61217}

Plato. (1588). \emph{Divini Platonis Opera Omnia}, edited by Ficino, M.,
Conrado, S. Lugduni: apud Joannem Lertout.
\url{https://doi.org/10.3931/e-rara-24349}

Plinius Secundus, G. (1250). \emph{Plinius Secundus Major, Historia
Naturalis. Liber Trigesimus Septimus Manu Recentiore Suppletus Est}.
Latin 6797. Paris: Bibliothèque nationale de France.
\url{https://gallica.bnf.fr/ark:/12148/btv1b550140045/}

Plotinus. (1580). \emph{Plotini Platonicorum facile coryphaei Operum
philosophicorum omnium libri LIV in sex Enneades distributi}. Edited by
Ficinus, M. Basileae: Ad Perneam Lecythum.
\url{https://doi.org/10.3931/e-rara-120889}

\end{document}
